\subsection{CHAMBERS AND BASES OF ROOT SYSTEMS}

For all \(\alpha \in \mathbb{R}\) , let \(L_{\alpha}\) be the hyperplane of \(V\) consisting of the points invariant under \(s_{\alpha}\) . The chambers in \(V\) determined by the set of the \(L_{\alpha}\) (Chap. V, § 1, no. 3) are called the chambers of \(R\) . The bijection \(V \to V^{*}\) defined by the scalar product \((x|y)\) takes \(\alpha\) to \(\frac{2\alpha^{-}}{(\alpha^{-}|\alpha^{-})}\) for \(\alpha \in R\) , hence \(L_{\alpha}\) to \(L_{\alpha^{-}}\) , and hence the chambers of \(R\) to those of \(R^{-}\) . If \(C\) is a chamber of \(R\) , the corresponding chamber of \(\mathbb{R}^{\sim}\) is denoted by \(\mathbb{C}^{\sim}\) . By Prop. 7 of no. 2, \(\mathbb{C}^{\sim}\) depends only on \(\mathbb{C}\) and not on the choice of \((x|y)\) .

THEOREM 2. (i) The group W(R) acts simply-transitively on the set of chambers.

\begin{enumerate}
\def\labelenumi{(\roman{enumi})}
\setcounter{enumi}{1}
\item
  Let \(\mathrm{C}\) be a chamber. Then \(\overline{\mathrm{C}}\) is a fundamental domain for \(\mathrm{W}(\mathrm{R})\) .
\item
  C is an open simplicial cone (Chap. V, § 1, no. 6).
\item
  Let \(L_{1}, L_{2}, \ldots, L_{l}\) be the walls of C. For all i, there exists a unique indivisible root \(\alpha_{i}\) such that \(L_{i} = L_{\alpha_{i}}\) and such that \(\alpha_{i}\) is on the same side of \(L_{i}\) as C.
\item
  The set B(C) = \(\{\alpha_{1},\ldots,\alpha_{l}\}\) is a basis of V.
\item
  C is the set of \(x \in V\) such that \(\langle \alpha_i, x \rangle > 0\) for all \(i\) (or, equivalently, the set of \(x \in V\) such that \((x|\alpha_i) > 0\) for all \(i\) ).
\item
  Let \(S\) be the set of the \(s_{\alpha_i}\) . The pair \((W(R), S)\) is a Coxeter system (Chap. IV, § 1, no. 3).
\end{enumerate}

Assertions (i) and (vii) follow from Chap. V, § 3, no. 2, Th. 1. Assertion (ii) follows from Chap. V, § 3, no. 3, Th. 2. Assertion (iv) is clear. The root \(\alpha_{i}\) is orthogonal to \(L_{i}\) , and \(\alpha_{i}^{-}\) is identified with \(2\alpha_{i}/(\alpha_{i}|\alpha_{i})\) . Since W(R) is essential (no. 1, Remark 3), assertions (iii), (v) and (vi) follow from Chap. V, § 3, no. 9, Prop. 7.

Remarks. 1) Assertion (vii) shows in particular that \(\mathrm{W}(\mathbb{R})\) is generated by the reflections \(s_{\alpha_i}\) .

\begin{enumerate}
\def\labelenumi{\arabic{enumi})}
\setcounter{enumi}{1}
\item
  If \(x, y \in \mathbb{C}\) , then \((x|y) > 0\) (Chap. V, § 3, no. 5, Lemma 6), in other words the angle \((\widehat{x,y})\) is acute.
\item
  Let \(m(\alpha, \beta)\) be the order of \(s_{\alpha} s_{\beta} (\alpha, \beta \in \mathrm{B}(\mathrm{C}))\) . The matrix \((m(\alpha, \beta))\) is identified with the Coxeter matrix (Chap. IV, § 1, no. 9) of (W, S). If \(\alpha \neq \beta\) , Prop. 3 of Chap. V, § 3, no. 4 shows that the angle \((\widehat{\alpha, \beta})\) is equal to \(\pi - \frac{\pi}{m(\alpha, \beta)}\) ; in particular, this angle is either obtuse or equal to \(\pi\) , and \((\alpha | \beta) \leqslant 0\) . By using the list in no. 3, it follows that \(m(\alpha, \beta)\) is equal to 2, 3, 4 or 6.
\end{enumerate}

DEFINITION 2. A subset B of R is called a basis of R if there exists a chamber C of R such that B = B(C). If C is a chamber, B(C) is called the basis of R defined by C.

Remarks. 4) Assertion (vi) of Th. 2 shows that the map C ↦ B(C) is a bijection from the set of chambers to the set of bases. Consequently, W(R) acts simply-transitively on the set of bases.

\begin{enumerate}
\def\labelenumi{\arabic{enumi})}
\setcounter{enumi}{4}
\tightlist
\item
  Let C be a chamber of R, and let B be the corresponding basis. If \(\alpha\in B\) , put \(\varphi(\alpha)=\alpha^{-}\) if \(2\alpha\notin R\) and \(\varphi(\alpha)=\frac{1}{2}\alpha^{-}\) if \(2\alpha\in R\) . Then \(\varphi(B)\) is the basis of \(R^{-}\) defined by \(C^{-}\) ; this follows from the fact that the walls of \(C^{-}\) are the \(L_{\alpha^{-}}\) , for \(\alpha\in B\) .
\end{enumerate}

DEFINITION 3. Let B be a basis of R. The Cartan matrix of R (relative to B) is the matrix \((n(\alpha,\beta))_{\alpha,\beta\in\mathbb{B}}\) .

For all \(\alpha \in \mathrm{B}\) , \(n(\alpha, \alpha) = 2\) . For \(\alpha, \beta \in \mathrm{B}\) ,

\[
n (\alpha , \beta) = 2 \frac {(\alpha | \beta)}{(\beta | \beta)} = - 2 \frac {\| \alpha \|}{\| \beta \|} \cos \frac {\pi}{m (\alpha , \beta)},\tag{11}
\]

where \(m(\alpha, \beta)\) denotes as above the order of \(s_{\alpha}s_{\beta}\) . If \(\alpha \neq \beta\) , \(n(\alpha, \beta) = 0\) , -1, -2 or -3 (cf.~no. 3).

Remarks. 6) The Cartan matrix \((n(\alpha,\beta))\) should not be confused with the Coxeter matrix \((m(\alpha,\beta))\) . Note in particular that the Cartan matrix is not necessarily symmetric.

\begin{enumerate}
\def\labelenumi{\arabic{enumi})}
\setcounter{enumi}{6}
\tightlist
\item
  Canonical indexing. If B and \(\mathbf{B}'\) are two bases of R, there exists a unique element \(w \in \mathbf{W}\) such that \(w(\mathbf{B}) = \mathbf{B}'\) . We have
\end{enumerate}

\[
n (w (\alpha), w (\beta)) = n (\alpha , \beta) \text { and } m (w (\alpha), w (\beta)) = m (\alpha , \beta)
\]

for \(\alpha, \beta \in \mathbf{B}\) . Consequently, the Cartan and Coxeter matrices associated to B can be obtained from those associated to \(\mathbf{B}'\) by composition with the bijection

\[
\alpha \mapsto w (\alpha)
\]

from B to B'.

The Cartan and Coxeter matrices can actually be defined canonically in the following way. Let X be the set of pairs \((B, \alpha)\) , where B is a basis of R and \(\alpha \in B\) . The group W acts in an obvious way on X and each orbit of W on X meets each of the sets \(\{B\} \times B\) in exactly one point. If I is the set of these orbits, each basis B admits a canonical indexing \((\alpha_{i})_{i \in I}\) . Moreover, there exists a unique matrix \(N = (n_{ij})\) (resp. \(M = (m_{ij})\) ), of type \(I \times I\) , such that for any basis B, the Cartan (resp. Coxeter) matrix associated to B can be obtained from N (resp. M) by composing with the canonical indexing of B; it is called the canonical Cartan matrix (resp. Coxeter matrix) of R.

PROPOSITION 15. Let B be a basis of R and \(\alpha\) an indivisible root. There exist \(\beta \in B\) and \(w \in W(R)\) such that \(\alpha = w(\beta)\) .

Let C be the chamber such that \(\mathrm{B} = \mathrm{B}(\mathrm{C})\) . The hyperplane \(L_{\alpha}\) is a wall of a chamber \(C'\) of R, and there exists an element of \(\mathrm{W}(\mathrm{R})\) that transforms \(C'\) to C. We are therefore reduced to the case where \(L_{\alpha}\) is a wall of C. Then \(\alpha\) is proportional to an element \(\beta\) of R. Since \(\alpha\) and \(\beta\) are indivisible, \(\alpha = \pm\beta\) . If \(\alpha = -\beta\) , then \(\alpha = s_{\beta}(\beta)\) , hence the proposition.

COROLLARY. Let \(R_{1}\) and \(R_{2}\) be two reduced root systems in vector spaces \(V_{1}\) and \(V_{2}\) , and let \(B_{1}\) and \(B_{2}\) be bases of \(R_{1}\) and \(R_{2}\) . Let \(f: B_{1} \to B_{2}\) be a bijection that transforms the Cartan matrix of \(R_{1}\) to that of \(R_{2}\) . Then there

exists an isomorphism \(\mathrm{F}: \mathrm{V}_1 \to \mathrm{V}_2\) that transforms \(\mathbb{R}_1\) to \(\mathbb{R}_2\) and \(\alpha\) to \(f(\alpha)\) for all \(\alpha \in \mathbb{B}_1\) .

Let F be the isomorphism from \(V_{1}\) to \(V_{2}\) that takes \(\alpha\) to \(f(\alpha)\) for all \(\alpha \in B_{1}\) . Then F transforms \(s_{\alpha}\) to \(s_{f(\alpha)}\) , hence \(W(R_{1})\) to \(W(R_{2})\) (Th. 2), and hence \(R_{1}\) to \(R_{2}\) (Prop. 15).

PROPOSITION 16. Let B be a basis of R, and G the subgroup of A(R) consisting of the elements leaving B stable. Then W(R) is a normal subgroup of A(R) and A(R) is the semi-direct product of G and W(R).

If \(\alpha\in R\) and \(t\in\mathrm{A}(\mathrm{R})\) , then \(ts_{\alpha}t^{-1}=s_{t(\alpha)}\) ; since \(\mathrm{W}(\mathrm{R})\) is generated by the \(s_{\alpha}\) , we see that \(\mathrm{W}(\mathrm{R})\) is a normal subgroup of \(\mathrm{A}(\mathrm{R})\) . By transport of structure, \(\mathrm{A}(\mathrm{R})\) transforms a basis of R to a basis of R. Since \(\mathrm{W}(\mathrm{R})\) acts simply-transitively on the set of bases, every element of \(\mathrm{A}(\mathrm{R})\) can be written uniquely in the form \(g_{1}g_{2}\) , where \(g_{1}\in\mathrm{W}(\mathrm{R})\) and \(g_{2}\in G\) .

Remarks. 8) Let \(\mathbb{R}_1, \ldots, \mathbb{R}_p\) be root systems in vector spaces \(\mathbb{V}_1, \ldots, \mathbb{V}_p\) , \(\mathbb{R}\) the direct sum of the \(\mathbb{R}_i\) in \(\mathbb{V} = \prod_i \mathbb{V}_i\) , \(\mathbb{C}_i\) a chamber of \(\mathbb{R}_i\) , and \(\mathbb{B}_i = \mathbb{B}(\mathbb{C}_i)\) . It is immediate that \(\mathbb{C} = \prod_i \mathbb{C}_i\) is a chamber of \(\mathbb{R}\) and that \(\mathbb{B}(\mathbb{C}) = \bigcup_i \mathbb{B}_i\) . It follows from Th. 2 that all the chambers and bases of \(\mathbb{R}\) are obtained in this way.

